% Part: incompleteness
% Chapter: theories-computability
% Section: inseparability

\documentclass[../../../include/open-logic-section]{subfiles}

\begin{document}

\olfileid{inc}{tcp}{ins}

\olsection{Ukara kang Bisa Dibuktekake lan Dibantah ing $\Th{Q}$ Ora Bisa
  Dipisahake kanthi Komputabel}


Upamana $\Th{\bar Q}$ iku himpunan ukara kang {\em negasine} bisa
dibuktekake ing $\Th{Q}$, yaiku $\Th{\bar Q} = \Setabs{!A}{\Th{Q} \Proves
  \lnot !A}$. Elinga yen himpunan $A$ lan $B$ kang ora irisan diarani
ora bisa dipisahake kanthi komputabel yen ora ana himpunan komputabel
$C$ saengga $A \subseteq C$ lan $B \subseteq \Complement{C}$.

\begin{lem}
$\Th{Q}$ lan $\Th{\bar Q}$ ora bisa dipisahake kanthi komputabel.
\end{lem}

\begin{proof}
Upamana $C$ himpunan komputabel saengga $\Th{Q} \subseteq C$ lan
$\Th{\bar Q} \subseteq \Complement{C}$. Upamana $R(x,y)$ iku relasi
\begin{quote}
$x$ ngodeake formula $!D(u)$ lan $!D(\num y)$ kalebu ing $C$.
\end{quote}
Kita bakal nuduhake yen $R(x,y)$ iku relasi komputabel universal, kang
nuwuhake kontradiksi.

Upamana $S(y)$ komputabel, direpresentasikake dening $!D_S(u)$ ing
$\Th{Q}$. Mula
\begin{eqnarray*}
S(\num n) & \lif & \Th{Q} \Proves !D_S(\num n) \\
& \lif & !D_S(\num n) \in C
\end{eqnarray*}
lan
\begin{eqnarray*}
\lnot S(\num n) & \lif & \Th{Q} \Proves \lnot !D_S(\num n) \\
& \lif & !D_S(\num n) \in \Th{\bar Q} \\
& \lif & !D_S(\num n) \not\in C
\end{eqnarray*}
Mula $S(y)$ ekuivalen karo $R(\#(!D_S(\num u)),y)$.
\end{proof}

\end{document}
